\section{Implemented Surface}

The serial verifier, \texttt{verify\_package}, first validates contests, then
walks \texttt{package.summaries}. For each summary it looks up the contest,
runs \texttt{verify\_contest\_selection\_sum}, and pushes an
\texttt{EquationPass} with equation id \texttt{contest\_selection\_sum}.
After that opening loop, it extends the report with status-event, batch,
lineage, reference, privacy, CVR-summary, audit, RLA, jurisdiction-adapter, and
manual-audit checks.

The parallel verifier, \texttt{verify\_package\_parallel}, keeps the same
contest validation and the same later extension sequence. Only the summary loop
is replaced. It calls \texttt{par\_iter()} over \texttt{package.summaries},
performs the same contest lookup and selection-sum equation check, collects the
result as \texttt{Result<Vec<EquationPass>, RcountCoreError>}, and then extends
the report with that vector before running the remaining serial checks.

\begin{center}
\small
\begin{tabularx}{\textwidth}{lX}
\toprule
Check & Failure caught or preserved \\
\midrule
\texttt{contest\_selection\_sum} & First parallelized family; edited summary
selection totals fail in both entry points with the same error. \\
\texttt{EquationPass} ordering & Parallel summary checks are collected before
the later verifier families are appended, preserving the public report order. \\
\texttt{verify\_package} & Baseline serial transcript and error behavior. \\
\texttt{verify\_package\_parallel} & Parallel entry point that must equal the
serial \texttt{VerificationReport} on accepted fixtures. \\
\bottomrule
\end{tabularx}
\end{center}

This is intentionally the smallest useful parallelization boundary. It exposes
Rayon in \texttt{rcount-core} through \texttt{rayon::prelude}, but it does not
move source-hash, batch-accounting, or audit checks into parallel execution.
Those checks remain part of the equivalence report because the parallel entry
point still runs them in serial after the collected selection-sum passes.
